Now that I have established how software systems can be automatically constructed from meaning using a semantic toolkit, in this section, I will talk about how this drives the future of debugging. 

In section \ref{subsec:shared-meaning}, I discussed how shared meaning leads to successful interactions between designs and how that can be used to construct a distributed system by leveraging semantically compatible interactions. In this section I will move beyond just the boundary between designs and apply that principle more generally and present the implications of semantic compatibility between all parts of a system.

When two humans are unable to communicate effectively to achieve their intentions, the root cause of failure can be identified as an incompatibility in meaning. If an apple means something different to me than it does to another human, there will be a misunderstanding in our communication about apples. If I ask for an apple and they give me a pencil, then there is a failure. If you were to debug the human interaction, the misalignment in meaning will predict the failure and an agreement on the meaning will resolve the failure.

The same principles apply to software systems and now that they are defined at the level of meaning using the semantic toolkit, they can be debugged through the established meaning prior to execution. If systems use the same meaning to construct every intention, then they are semantically compatible by definition. So the system can be fully debugged by constructing it with shared meaning and reasoning about that meaning.

This is not to say that the program can't fail, the shared meaning itself can be incomplete and through failures, reality will reveal the missing semantics after root cause analysis. However, after learning the new semantics, the meaning of the semantic toolkit is updated and this means that the intentions of all parts of the system are realized from the same meaning, ensuring semantic compatibility. Then through reasoning about the meaning, the engine can resolve all the known failures in the closed semantic world.

This is much more precise than human communication. In human communication, the participants work to ensure that they have the same meaning and human skill plays a part in how effectively the meaning can be remembered and applied. However, in machine communication, they literally share the same meaning. A great analogy would be that it is the same mind realizing its intentions in multiple nodes, it can't misunderstand itself because it is communicating with itself. 

\begin{figure}[htbp]
    \centering
    \includesvg[
        width=0.90\linewidth
    ]{assets/future_debugging.svg}
    \caption{Eliminating bugs within and between systems using semantic toolkit}
    \label{fig:future_debugging}
\end{figure}

The future of debugging is ensuring that systems understand each other through shared meaning by constructing the meaning of their intentions from the same semantic toolkit and reasoning about their semantics to ensure that all known failures in the closed semantic world are eliminated. This process eliminates misunderstanding between and within systems and eliminates failures before they manifest.

Ultimately, this has transformed debugging from a reactive task after a failure happens to a proactive task that eliminates bugs by construction. As a result, systems will not fail in known ways. If they do fail an automatic process will collect the necessary diagnostic information to learn new semantics through root cause analysis. The recurrence of the same failure and/or potential new failures are eliminated by sharing the updated meaning across all nodes in the system and allowing the engine to reason about them.

In the next section, in addition to debugging through shared meaning, I will describe how systems of the future will be deterministically intelligent by reasoning about the meaning through their narratives.